% ED26-REQ-N-c-004-03: Immerman 1986, selected sections 1--3.

\subsection{Immerman 1986: ordered polynomial-time queries}
\label{sec:p09-n-c-004-03}

Immerman's ordered-capture theorem describes a \emph{fixed} relational
query, not the complexity of taking arbitrary formulas as input.  Its hard
direction constructs a single least fixed point describing a deterministic
machine's successive tape rows.  The construction below retains the
article's initial-row formula, three-cell transition table, and final
accepting-cell readout \cite[secs.~1--3, thm.~1 and lem.~3.1]{Immerman1986}.
The inspected author-hosted copy identifies the 1986 journal article but
has its own pagination, 1--18; its selected definitions and proof are on
author-copy pp.~2--8, corresponding to the journal locators pp.~87--93.
It is not the six-page 1982 preliminary presentation or the 1999 book.

\begin{hbdefinition}{def:p09-immerman1986-query}{Databases, queries, and the source's input convention}
A similarity type
\(\tau=\langle R_1^{a_1},\ldots,R_j^{a_j},c_0,\ldots,c_s\rangle\)
is a fixed finite list of relation and constant symbols.  A database
\(B=\langle U,R_1^B,\ldots,R_j^B,c_0^B,\ldots,c_s^B\rangle\)
is a finite nonempty structure of this type, with \(R_i^B\subseteq U^{a_i}\)
and \(c_i^B\in U\).  Put \(n=\lvert U\rvert\).
The source's domain relational calculus \(\mathcal L(\tau)\) is
first-order logic with equality, these symbols, Boolean operations and
quantification over \(U\).  A formula \(\gamma(\bar x)\), with a
fixed \(d\)-tuple of free variables, returns
\[
 \gamma^B=\{\bar a\in U^d:B,\bar x\mapsto\bar a\entails\gamma\}.
\]
A sentence is a Boolean query.  Query membership means deciding whether
\(\bar a\in\gamma^B\); computing the entire answer relation is
also polynomially related when \(d\) is fixed.
Queries respect isomorphisms of their input type: an isomorphism maps
the answer tuples to the corresponding answer tuples.  In the ordered
case the distinguished order is one of the symbols it must preserve.

The article implicitly encodes a database by a binary string listing the
characteristic functions of its relations.  For fixed \(\tau\), dense
relation tables, explicit domain size, and representations of its finitely
many constants have size polynomial in \(n\).  Constants can equivalently
be recorded by their singleton unary tables.  These conventions are fixed
before discussing a machine.  Write \(\mathrm{QPTIME}\) for fixed
queries whose membership graphs
\(\{(B,\bar a):\bar a\in Q(B)\}\) are decidable in time
polynomial in \(n\), under such explicit encodings.  No exponentially
succinct database representation is included.
\end{hbdefinition}

\begin{hbdefinition}{def:p09-immerman1986-lfp}{The positive fixed-point language}
For a formula \(\varphi(R,\bar x,\bar z)\), an \(r\)-ary relation
variable \(R\), an \(r\)-tuple of distinct variables \(\bar x\),
and a fixed assignment \(\bar b\) to the remaining parameters, define
\[
 \Gamma_{\varphi,\bar b}^B(S)
 =\{\bar a\in U^r:(B,R\mapsto S),
                  (\bar x,\bar z)\mapsto(\bar a,\bar b)\entails\varphi\}.
\]
For the positive-arity binder convention underlying the source's
constant proviso, \(r\geq1\); the output query may still be Boolean.
The binder is admitted only when \(R\) occurs positively: with the
source's primitive negation and conjunction syntax, its occurrences lie
under an even number of negations.  Bound relation variables and ordinary
variables are renamed to prevent capture.  Positivity makes this operator
monotone.  Its least fixed point is computed from
\[
 I_0=\varnothing,\qquad I_{i+1}=\Gamma_{\varphi,\bar b}^B(I_i).
\]
The source writes \((\bar t:\mu R(\bar x))\varphi\), abbreviated
\((\mu R)\varphi\), for membership of the value of \(\bar t\)
in that least fixed point.  Its language \(\mathcal L_\mu(\tau)\)
closes \(\mathcal L(\tau)\) under this operation and ordinary Boolean
operations and quantification, so nested fixed points are allowed.
\end{hbdefinition}

\begin{hblemma}{thm:p09-immerman1986-evaluation}{proved=annex:p9.n-c-004.3.evaluation}{Least fixed points and fixed-query polynomial time}
For every fixed \(\mathcal L_\mu\) query \(\gamma\), there is a
polynomial \(p_\gamma\) such that \(\gamma^B\) is computable in
time \(p_\gamma(n)\).  An \(r\)-ary monotone induction from the
empty relation stabilizes after at most \(n^r\) strict increases.
This is the source's Fact~2.1, including the requirement that the query
and its relation arities be fixed \cite[sec.~2]{Immerman1986}.
\end{hblemma}

\begin{hbexample}{def:p09-immerman1986-family}{Grandparent and the actual ancestor operator}
The source's family type has a unary female predicate \(F\), binary
parent and husband predicates \(P,H\), and one constant, for example
Abraham.  Grandparent is the first-order query
\(\exists z(P(x,z)\land P(z,y))\).  Its fixed-point example is
\[
 \varphi(R)[x,y]\defeq x=y\lor\exists z(P(x,z)\land R(z,y)).
\]
The stages start with the empty relation, then the diagonal, then the
diagonal together with parent pairs, then paths of length at most two,
and so forth.  Thus the displayed operator defines the \emph{reflexive}
transitive closure of \(P\): it includes paths of length zero.
The source calls this ancestor; removing the diagonal term would instead
require a different base clause for positive-length ancestry.  For a
chain Abraham--Isaac--Jacob, its first three nonempty stages are the
diagonal, the diagonal plus the two parent pairs, and those pairs plus
(Abraham, Jacob).  Husband and female data do not enter this operator.
\end{hbexample}

\subsubsection{Order, tuple coordinates, and the first tape row}

In the capture theorem, \(\leq\) is an available total ordering on
\(U\); write \(x<y\) for \(x\leq y\land x\ne y\).
For each fixed \(k\), lexicographic order on \(U^k\) is first-order:
\[
 \operatorname{Lex}_k(\bar u,\bar v)\defeq
 \bigvee_{i=1}^k\left(\bigwedge_{j<i}u_j=v_j\land u_i<v_i\right).
\]
Its successor is
\[
 \operatorname{Succ}_k(\bar u,\bar v)\defeq
 \operatorname{Lex}_k(\bar u,\bar v)\land
 \neg\exists\bar w(
   \operatorname{Lex}_k(\bar u,\bar w)\land
   \operatorname{Lex}_k(\bar w,\bar v)).
\]
First and last tuple predicates are defined by the absence of a smaller
or larger tuple.  These finite formulas supply the article's abbreviations
\(\bar p-1,\bar p+1,\bar t-1\); they do not posit arithmetic
functions on the database.  Identifying the ordered domain with
\(0,\ldots,n-1\), \(k\)-tuples address \(n^k\) tape positions
or times.  The two coordinate orders may be different, as specified below.

\begin{hblemma}{thm:p09-immerman1986-initial}{proved=annex:p9.n-c-004.3.initial}{The source initial-row formula and general-type decoding}
There is a first-order formula \(M_0(\bar p,a)\) defining the initial
letter at each tape position for a fixed polynomial-time machine and a
fixed database type, once the finite alphabet is coded by distinct domain
elements.  The finitely many domains too small to provide the codes are
handled separately, not assigned nonexistent code elements.

The source prints the case of one binary input relation \(E\) and
\(k=3\): the first \(n^2\) positions list its characteristic bits,
and the remaining \(n^3-n^2\) positions are blank.  Here
\(0,1\) code bits, \(2\) blank, and \(3,4\) the start state
looking at bit zero or one.  With these numerals denoting the corresponding
ordered code elements,
\[
\begin{split}
 M_0(p_1,p_2,p_3,a)\defeq{}&
 (p_1=p_2=p_3=0)\land
  ((E(0,0)\land a=4)\lor(\neg E(0,0)\land a=3))\\
 &{}\lor(p_1=0\land(p_2\ne0\lor p_3\ne0)
       \land(a=0\lor a=1)
       \land(E(p_2,p_3)\leftrightarrow a=1))\\
 &{}\lor(p_1\ne0\land a=2).
\end{split}
\]
No fixed-point relation is used in this initial-row definition.
For arbitrary fixed types the annex expands the source's
encoding/decoding generalization by order-definable characteristic-table
blocks, including constants; the printed binary layout above is retained
as its own case \cite[lem.~3.1]{Immerman1986}.
\end{hblemma}

\subsubsection{The local transition relation and the accepting fixed point}

Normalize a fixed deterministic polynomial-time machine \(M\) to a
one-tape machine whose head moves by at most one cell.  Its finite alphabet
includes ordinary tape letters and letters carrying a state/head marker.
On acceptance it returns to position zero, writes a designated accepting
letter \(q_f\) there, and thereafter leaves the configuration unchanged;
a distinct rejecting letter may be used on rejection.  One-way left-end
behavior is fixed.  These changes cost only polynomial time.  Choose
\(k\) and a fixed small-domain cutoff so that on larger domains this
normalized computation, including input decoding and the return, ends
before \(n^k-1\), never visits the last tape cell, and has room for
the fixed alphabet codes.  The final tableau can then be padded with
stationary halting rows.

The source writes \(\langle xyz\rangle\to w\) for a finite local
transition table \(\delta_M\).  Its first-order expansion is
\[
 \Delta_M(x,y,z,w)\defeq
 \bigvee_{(b_{-1},b_0,b_1,b)\in\delta_M}
 (x=b_{-1}\land y=b_0\land z=b_1\land w=b).
\]
The entries are letter codes, not names of database constants.  Every
letter equality \(v=b\) in these displays abbreviates the first-order
test \(\operatorname{Code}_b(v)\) for its prescribed fixed-rank code
element (or \(v=c_0\) for the anchored accepting code).  Expand
these tests in the operator; they remain well-formed formulas even on
small domains where some code tests are false.  No partial code term
or new input constant is introduced.  The
matrix inspects the previous letters at the cell immediately to the left,
the cell itself, and the cell immediately to the right.  An absent neighbor
at a tape endpoint is a specified exterior letter \(\#\), not an
undefined predecessor or a wraparound cell.

For an induction relation \(R(\bar p,\bar t,a)\) of arity \(2k+1\),
let \(L_R(\bar p,\bar t,a)\) be
\[
 (\operatorname{First}_k^p(\bar p)\land a=\#)
 \lor\exists\bar u(
    \operatorname{Succ}_k^p(\bar u,\bar p)\land R(\bar u,\bar t,a)),
\]
and let \(H_R\) be its right-hand version with \(\operatorname{Last}\)
and \(\operatorname{Succ}_k^p(\bar p,\bar u)\).
Superscripts \(p,t\) specify the position or time order.  Expand the
source's tableau operator, with its endpoint clauses, as
\[
\begin{split}
 \varphi_M(R)[\bar p,\bar t,a]\defeq{}&
 (\operatorname{First}_k^t(\bar t)\land M_0(\bar p,a))\\
 &{}\lor\exists\bar u\exists a_{-1}a_0a_1\bigl[
 \operatorname{Succ}_k^t(\bar u,\bar t)
 \land\Delta_M(a_{-1},a_0,a_1,a)\\
 &\hspace{36mm}{}\land L_R(\bar p,\bar u,a_{-1})
 \land R(\bar p,\bar u,a_0)
 \land H_R(\bar p,\bar u,a_1)\bigr].
\end{split}
\]
All occurrences of \(R\) are positive.  Tape bit zero is itself
represented by a positive \(R\)-fact with letter code zero; the
formula does not attempt to infer a zero bit from a missing \(R\)-fact.

\begin{hblemma}{thm:p09-immerman1986-tableau}{proved=annex:p9.n-c-004.3.tableau}{The positive induction generates exactly the machine rows}
Let \(C_M=(\mu R)\varphi_M\).  In a database large enough for the
chosen codes, \(C_M(\bar p,\bar t,a)\) holds exactly when the
machine's tape letter at that position and time is \(a\).
Iteration zero is empty; iteration \(i\), for \(1\leq i\leq n^k\),
contains precisely the first \(i\) complete rows and no later row.
Each position in each produced row has exactly one letter.  Consequently
the accepting readout is the source's
\(C_M(\bar0,\overline{n^k-1},q_f)\), using the indicated first
position and last time, rather than merely the existence of an arbitrary
partial accepting configuration.
\end{hblemma}

\begin{hbtheorem}{thm:p09-immerman1986-capture}{proved=annex:p9.n-c-004.3.capture}{Theorem 1: ordered capture in the article's syntax}
On finite databases with a distinguished total order, and with the
source's at-least-one-constant convention,
\[
 \Sigma_\omega(\leq)=\mathrm{QPTIME}=\mathcal L_\mu(\leq).
\]
Here \(\Sigma_\omega\) is the article's one-fixed-point level:
a single outer \(\mu\) with a first-order defining formula.  It is
not an arithmetical-hierarchy class.  The formula and vocabulary are fixed
while the input database varies.  The source proves the construction for
Boolean queries for simplicity; the annex gives the return to fixed-arity
query relations by coding the tested output tuple as parameters.

For a Boolean query, its constant \(c_0\) permits the single fixed-point
expression \(C_M(c_0,\ldots,c_0)\): the position order puts \(c_0\)
first, the time order puts \(c_0\) last, and the alphabet coding makes
\(c_0\) the accepting letter.  These are three different typed uses of
one element, as in the source's footnote~3.  Finite small-domain exceptions
are included by complete first-order diagrams and can be folded into the
same positive operator.  Without a constant, the corresponding normal
form allows one final ordinary quantifier, \(\exists z\), after the
fixed-point expression, with \(z\) serving as its coding anchor.
\end{hbtheorem}

The order is not a disposable label: it supplies first-order tape and time
successors and characteristic-table coordinates.  Queries on ordered
structures respect isomorphisms of \emph{all} their symbols, including
the order; forgetting that order changes the input class.  The source's
separate Theorem~2 concerns collapse of the unordered fixed-point hierarchy
and has its own section~4 proof.  Neither that collapse nor the 1986
capture proof supplies the separately required 1999 Chapter~4 treatment.

\subsubsection{Input axes and the unordered boundary}

\begin{table}[htbp]
\centering\small
\begin{tabular}{@{}p{0.21\textwidth}p{0.28\textwidth}p{0.43\textwidth}@{}}
\toprule
Task & Variable input; fixed data & Conclusion here \\
\midrule
Fixed-query evaluation & Database \(B\) and tested tuple vary;
query, type, arities and output arity fixed & Polynomial time by
Fact~2.1; output enumeration has a fixed polynomial factor. \\
Expression complexity & Formula varies; database fixed & Not the
input axis bounded by Fact~2.1. Formula length, variable count,
fixed-point arity and nesting are no longer constants. \\
Combined evaluation & Both database and formula vary & No uniform
polynomial-time bound follows from the fixed-query proof. \\
Ordered descriptive capture & One fixed query/property on ordered
databases; its membership graph in polynomial time & Theorem~1 constructs
one fixed positive induction for that query. It is a class correspondence,
not a combined model-checking bound or a test for whether an arbitrary
program runs in polynomial time. \\
Unordered parity & Database varies without an order;
formula and type fixed & Polynomial-time counting query, but not
\(\mathcal L_\mu\)-definable; order cannot be deleted from capture. \\
\bottomrule
\end{tabular}
\caption{Task and input axes for the 1986 capture statement.}
\label{tab:p09-immerman1986-input-axes}
\end{table}

\begin{hbtheorem}{thm:p09-immerman1986-unordered}{proved=annex:p9.n-c-004.3.unordered}{The source's even-number-of-females boundary}
In the unordered family type \(\langle F,P,H,c_0\rangle\), the
Boolean query ``\(\lvert F\rvert\) is even'' is polynomial-time
computable but not expressible in \(\mathcal L_\mu\).
For every fixed sentence, its truth is eventually constant on the databases
with \(F=U\), \(P=H=\varnothing\), and one distinguished \(c_0\).
Parity is not eventually constant on that subclass.  This is the example
reported in Fact~2.2, attributed in the article to Chandra and Harel;
the annex supplies the explicit equality-pattern argument for this example
\cite[sec.~2, fact~2.2]{Immerman1986}.
\end{hbtheorem}

\begin{proofboundary}{def:p09-immerman1986-standing}{Source construction and editorial completion}
The source prints the binary initial-row case and the positive
three-neighbor operator.  Its general-type decoding, machine normalization,
endpoint values, row-by-row exactness, and finite-exception integration are
expanded here as explicit completions of that construction.  The
article's constant coding is retained, including its alternative final
quantifier when no constant is available.  No 1982, 1999, or generic
successor-tableau account is used in place of its \(2k+1\)-ary
cell/time/letter induction.  The independently paginated author copy is
qualified as that copy, not misidentified as a journal facsimile.
\end{proofboundary}
